Nach dem Stefan-Boltzmann-Gesetz ist die Strahlungsleistung $P = \varepsilon\,\sigma\,A\,T^4$ proportional zu $T^4$. Aus $P_2/P_1 = \fPO$ folgt
\begin{align}
 \frac{T_2}{T_1} &= \fTF = \sqrt[4]{\fPO} = \fT \approx \result{\fTP{0}{3}}
\end{align}
Die \emph{absolute} Temperatur muss also um $\pTP{0}{2}$ zunehmen.
